\documentclass[12pt]{article}
\usepackage{amsmath, amssymb, amsthm}
\usepackage{geometry}
\usepackage{enumitem}
\usepackage{graphicx}

\geometry{margin=1in}

\title{Lecture Scribe}
\author{Course: CSE400 -- Fundamentals of Probability in Computing}
\date{Lecture 6 \\ Topic: Discrete Random Variables, Expectation, and Problem Solving \\ January 22, 2025}

\begin{document}
\maketitle

\section*{Definitions}

\subsection*{Random Variable}

A random variable $X$ on a sample space $\Omega$ is a function
\[
X : \Omega \rightarrow \mathbb{R}
\]
that assigns a real number to each outcome $\omega \in \Omega$.

In this lecture, we restrict attention to discrete random variables, i.e., random variables that take values from a finite or countably infinite set.

\subsection*{Discrete Random Variable}

A random variable is discrete if:
\begin{itemize}
    \item Its support is countable (finite or countably infinite),
    \item Probabilities are assigned to individual values,
    \item Each possible value has strictly positive probability.
\end{itemize}

\subsection*{Probability Mass Function (PMF)}

Let $X$ be a discrete random variable with range
\[
R_X = \{x_1, x_2, x_3, \dots\}.
\]

The probability mass function (PMF) of $X$ is defined as:
\[
P_X(x_k) = P(X = x_k), \quad k = 1,2,3,\dots
\]

\section*{Theorems / Propositions}

\subsection*{Normalization Property of PMF}

Since a discrete random variable must take exactly one of its possible values,
\[
\sum_{k=1}^{\infty} P_X(x_k) = 1.
\]

\subsection*{Bayes' Formula (Proposition 3.1)}

Given events $B_1, B_2, \dots, B_n$ forming a partition of the sample space and an event $A$ with $P(A) > 0$,
\[
P(B_i \mid A) =
\frac{P(A \mid B_i)P(B_i)}
{\sum_{j=1}^{n} P(A \mid B_j)P(B_j)}.
\]

\begin{itemize}
    \item $P(B_i)$: a priori probability
    \item $P(B_i \mid A)$: posterior probability
\end{itemize}

\section*{Proofs}

\subsection*{Proof of PMF Normalization}

Since $X$ must take one of the values $x_k$,
\[
1 = P\left(\bigcup_{k=1}^{\infty} \{X = x_k\}\right).
\]

Because the events $\{X = x_k\}$ are mutually exclusive,
\[
1 = \sum_{k=1}^{\infty} P(X = x_k) = \sum_{k=1}^{\infty} P_X(x_k).
\]

\subsection*{Derivation of Bayes' Formula}

Using conditional probability,
\[
P(A \mid B_i) = \frac{P(A \cap B_i)}{P(B_i)}
\Rightarrow
P(A \cap B_i) = P(A \mid B_i)P(B_i).
\]

Since
\[
A = \bigcup_{j=1}^{n} (A \cap B_j),
\]
we have
\[
P(A) = \sum_{j=1}^{n} P(A \mid B_j)P(B_j).
\]

Substituting into
\[
P(B_i \mid A) = \frac{P(A \cap B_i)}{P(A)}
\]
gives Bayes' formula.

\section*{Derivations / Equations}

\subsection*{Poisson-Type PMF Normalization}

Given
\[
p(i) = c \frac{\lambda^i}{i!}, \quad i = 0,1,2,\dots, \ \lambda > 0,
\]
using normalization,
\[
\sum_{i=0}^{\infty} p(i) = 1
\Rightarrow
c \sum_{i=0}^{\infty} \frac{\lambda^i}{i!} = 1.
\]

Since
\[
\sum_{i=0}^{\infty} \frac{\lambda^i}{i!} = e^{\lambda},
\]
we get
\[
c e^{\lambda} = 1
\Rightarrow
c = e^{-\lambda}.
\]

\section*{Examples}

\subsection*{Example 1: Tossing Three Fair Coins}

An experiment consists of tossing 3 fair coins.  
Let $Y$ denote the number of heads observed.

\textbf{Possible values:}
\[
Y \in \{0,1,2,3\}.
\]

\[
P(Y=0) = P(T,T,T) = \frac{1}{8}
\]
\[
P(Y=1) = \frac{3}{8}
\]
\[
P(Y=2) = \frac{3}{8}
\]
\[
P(Y=3) = \frac{1}{8}
\]

\textbf{Check normalization:}
\[
\sum_{i=0}^{3} P(Y=i) = 1.
\]

\subsection*{Example 2: PMF with Unknown Constant}

Given
\[
p(i) = c \frac{\lambda^i}{i!}, \quad i = 0,1,2,\dots
\]

From normalization,
\[
c = e^{-\lambda}.
\]

\[
P(X=0) = e^{-\lambda}.
\]

\[
P(X > 2) = 1 - P(X \leq 2)
\]
\[
= 1 - [P(X=0) + P(X=1) + P(X=2)]
\]
\[
= 1 - \left(e^{-\lambda} + \lambda e^{-\lambda} + \frac{\lambda^2}{2} e^{-\lambda}\right).
\]

\subsection*{Example 3: Bayes' Theorem -- Auditorium Problem}

There are 30 rows in an auditorium.  
Row $i$ has $10+i$ seats.

A row is chosen uniformly at random.  
A seat is chosen uniformly from the selected row.

\[
P(\text{Seat 15} \mid \text{Row 20}) = \frac{1}{30}.
\]

Using Bayes' theorem,
\[
P(\text{Row 20} \mid \text{Seat 15}) =
\frac{P(\text{Seat 15} \mid \text{Row 20})P(\text{Row 20})}
{\sum_{k=5}^{30} P(\text{Seat 15} \mid \text{Row } k)P(\text{Row } k)}
\]

\[
=
\frac{\frac{1}{30} \cdot \frac{1}{30}}
{\sum_{k=5}^{30} \frac{1}{10+k} \cdot \frac{1}{30}}
\approx 0.0325.
\]

\section*{Remarks}

\begin{itemize}
    \item PMF applies only to discrete random variables.
    \item The sum of PMF values must always equal 1.
    \item Bayes' theorem updates prior beliefs into posterior beliefs.
\end{itemize}

\section*{Assumptions}

\begin{itemize}
    \item Outcomes are equally likely unless stated otherwise.
    \item Coin tosses and selections are independent.
    \item All rows are selected with equal probability.
\end{itemize}

\section*{Special Cases}

\begin{itemize}
    \item Single-value random variable $\Rightarrow$ PMF equals 1 at that value.
    \item Continuous random variables use PDFs instead of PMFs.
    \item Individual points have zero probability in continuous random variables.
\end{itemize}

\section*{Summary: Key Results}

\begin{itemize}
    \item A random variable is a function from outcomes to real numbers.
    \item Discrete random variables use PMFs.
    \item $\sum P_X(x) = 1$.
    \item Bayes' theorem relates prior and posterior probabilities.
    \item Poisson-type PMFs normalize using $e^{\lambda}$.
\end{itemize}

\end{document}
